\begin{figure}
\centering\includegraphics[scale=0.7]{PyMPS_InfiniteDMRG_MPS.eps}
\caption{MPS structure generated by infinite-system DMRG at step $\ell$: a string of left-normalized $A$, a string of right-normalized $B$, joined by a diagonal singular value matrix $\Lambda^{[\ell]}$. Note that structurally the central unit does not repeat.}
\label{fig:infiniteDMRG_AB}
\end{figure}

Of course, at each step we discard the smallest singular values and their associated singular vectors once matrix dimensions exceed $D$, which is nothing but the density-matrix based truncation in the original formulation.
At each step (new chain length) we can write down $\hat{H}$ for that length as an MPO and do the energy minimization. Other operators find similar representations as in the finite-size case. 

Let me briefly go through the reformulation of this algorithm in the $\Gamma\Lambda$-notation. In the first step we simply rename $A$ and $B$ into $\Gamma$, in line with the translation of boundary sites in the finite-system case:
\begin{equation}
\ket{\psi_1} =  \sum_{\sigma_1^A \sigma_1^B} \Gamma^{[1]\sigma_1^A} \Lambda^{[1]} \Gamma^{[1]\sigma_1^B} \ket{\sigma_1^A \sigma_1^B} .
\end{equation}
We then minimize $\Psi^{\sigma_2^A \sigma_2^B}$ in 
\begin{equation}
\ket{\psi_2} = \sum_{\sigma_1^A \sigma_2^A \sigma_2^B \sigma_1^B} \Gamma^{[1]\sigma_1^A} 
\Psi^{\sigma_2^A \sigma_2^B} 
\Gamma^{[1]\sigma_1^B} \ket{\sigma_1^A \sigma_2^A \sigma_2^B \sigma_1^B},
\end{equation}
and decompose it -- as before -- into $A^{[2]\sigma_2^A} \Lambda^{[2]} B^{[2]\sigma_2^B}$. Now we define (and due to the labelling, there is a slight change for the $B$-matrices compared to the finite-system setup)
\begin{equation}
\Lambda^{[1]}_a \Gamma^{\sigma_2^A}_{ab} = A^{[2]\sigma_2^A}_{ab} \quad\quad
\Gamma^{\sigma_2^B}_{ab} \Lambda^{[1]}_b = B^{[2]\sigma_2^B}_{ab}
\end{equation}
to arrive at
\begin{equation}
\ket{\psi_2} =  \sum_{\sigma_1^A \sigma_2^A \sigma_1^B \sigma_2^B} \Gamma^{[1]\sigma_1^A} \Lambda^{[1]} \Gamma^{[2]\sigma_2^A} \Lambda^{[2]} \Gamma^{[2]\sigma_2^B} \Lambda^{[1]} \Gamma^{[1]\sigma_1^B}   \ket{\sigma_1^A \sigma_2^A \sigma_2^B \sigma_1^B} ,
\end{equation}
as represented in Fig.~\ref{fig:infiniteDMRG_GL}. 

We can now ask two questions: (i) in DMRG, finding the ground state by an iterative solver like Lanczos is the most time-consuming part. Can we find a speed-up by providing a good initial guess? In finite-system DMRG the MPS formulation automatically yielded White's prediction method, whereas attempts at speeding up infinite-system DMRG have been made in the past, meeting with mixed success\cite{Schollwoeck98,Sun02}. (ii) Can we use the information at the chain center to build a translationally invariant state (up to period 2) in the thermodynamic limit, find its ground state or evolve it in time?

The answer is yes to both questions, and builds on the identification of a two-site repetititve structure in the states.  As opposed to the $A,B$-notation, where the central unit does not repeat itself even in the thermodynamic limit, it is very easy to read off a two-site repeat unit in the $\Gamma\Lambda$-notation, given by
\begin{equation}
\Lambda^{[\ell-1]} \Gamma^{[\ell]\sigma_\ell^A} \Lambda^{[\ell]} \Gamma^{[\ell]\sigma_\ell^B} .
\end{equation}
Using the translation rules it can be translated into the $A$, $B$-language:
\begin{equation}
A^{[\ell]\sigma_\ell^A} \Lambda^{[\ell]} B^{[\ell]\sigma_\ell^B} (\Lambda^{[\ell-1]})^{-1} .
\end{equation}
This result can also be obtained directly from the $A$, $B$ notation, but the argument is more complicated than in the $\Gamma\Lambda$ notation. It is of course to be understood that repeating these state fragments does not generate the state they were taken from; the claim is just that in the thermodynamic limit $\ell\rightarrow\infty$, when all sites are created equal, this repetition can come close. In any case, they are an educated guess about what the state will look like! 

\begin{figure}
\centering\includegraphics[scale=0.7]{PyMPS_InfiniteDMRG_TEBD.eps}
\caption{MPS structure generated by infinite-system DMRG at step $\ell$ in the $\Gamma\Lambda$-notation: $\Gamma$ and $\Lambda$ matrices alternate, and a possible identification of a (repetitive) two-site building block is given.}
\label{fig:infiniteDMRG_GL}
\end{figure}

We will now put this state fragment to multiple use, first on finite systems generated by infinite-system DMRG and then on thermodynamic limit states, both in the context of ground state searches and time evolutions. In the former case, it will provide a highly efficient guess for the next quantum state; the evaluation of observables on this state proceed exactly as in the other finite systems.
In the second case, both ground state and time evolution algorithms can be formulated (iDMRG and iTEBD), which however necessitate both an (identical) analysis of the issue of orthonormality of states in the thermodynamic limit.

\subsection{State prediction in infinite-system DMRG} 
The identification of the ``unit cell'' of the state allows us to define a good initial guess for infinite-system DMRG\cite{McCulloch08}, which avoids all the problems encountered by previous authors and leads to a dramatic speed-up even for small chains, where the underlying assumption that the chain center is representative of the physics of the thermodynamic limit state is certainly wrong: in order to grow the chain, we simply insert one unit cell, even though for small chains the idea that the state is just a repetition of these unit cells is not well verified -- but even then so much better than a random guess. Starting from
\begin{equation}
\ket{\psi_{\ell}} =  \sum_{\fat{\sigma}} A^{[1]\sigma_1^A} \ldots A^{[\ell-1]\sigma_{\ell-1}^A} (A^{[\ell]\sigma_\ell^A} \Lambda^{[\ell]} B^{[\ell]\sigma_\ell^B} [\Lambda^{[\ell-1]}]^{-1}) \Lambda^{[\ell-1]} B^{[\ell-1]\sigma_{\ell-1}^B} \ldots B^{[1]\sigma_1^B} \ket{\fat{\sigma}},
\end{equation} 
where the repeat unit has been bracketed out, the guess will then read
\begin{eqnarray}
\ket{\psi_{\ell+1}^{{\rm guess}}} &=& \sum_{\fat{\sigma}} A^{[1]\sigma_1^A} \ldots  A^{[\ell-1]\sigma_{\ell-1}^A} \times \nonumber \\
& & ( A^{[\ell]\sigma_\ell^A} \Lambda^{[\ell]} B^{[\ell]\sigma_{\ell+1}^A} [\Lambda^{[\ell-1]}]^{-1} ) (A^{[\ell]\sigma_{\ell+1}^B} \Lambda^{[\ell]} B^{[\ell]\sigma_\ell^B}  
[\Lambda^{[\ell-1]}]^{-1} ) \times \nonumber \\
& &  \Lambda^{[\ell-1]} B^{[\ell-1]\sigma_{\ell-1}^B} \ldots B^{[1]\sigma_1^B} \ket{\fat{\sigma}} 
\end{eqnarray}
or, multiplying out, 
\begin{equation}
\ket{\psi_{\ell+1}^{{\rm guess}}} = \sum_{\fat{\sigma}} A^{[1]\sigma_1^A} \ldots A^{[\ell]\sigma_\ell^A} \Lambda^{[\ell]} B^{[\ell]\sigma_{\ell+1}^A} [\Lambda^{[\ell-1]}]^{-1} A^{[\ell]\sigma_{\ell+1}^B} \Lambda^{[\ell]} B^{[\ell]\sigma_\ell^B} B^{[\ell-1]\sigma_{\ell-1}^B} \ldots B^{[1]\sigma_1^B} \ket{\fat{\sigma}} .
\end{equation}
In this ansatz, we can now identify a guess for $\Psi^{\sigma_{\ell+1}^A\sigma_{\ell+1}^B}$ as
\begin{equation}
\Psi^{\sigma_{\ell+1}^A\sigma_{\ell+1}^B}_{{\rm guess}} =
\Lambda^{[\ell]} B^{[\ell]\sigma_{\ell+1}^A} [\Lambda^{[\ell-1]}]^{-1} A^{[\ell]\sigma_{\ell+1}^B} \Lambda^{[\ell]} .
\label{eq:guess2sites}
\end{equation}
From this ansatz, we can then iteratively find the $\Psi^{\sigma_{\ell+1}^A\sigma_{\ell+1}^B}$ that minimizes the energy in the infinite-system DMRG framework, generating from it $A^{[\ell+1]\sigma_{\ell+1}^A}$, $\Lambda^{[\ell+1]}$, and $B^{[\ell+1]\sigma_{\ell+1}^B}$.

Alternatively, the ansatz can be brought into in a more elegant form. At the moment, $B$-matrices show up on the A-side of the lattice and vice versa. But we can exploit our ability to canonize MPS, and canonize $\Lambda^{[\ell]} B^{[\ell]\sigma_{\ell+1}^A}$ by SVD to
$A^{[\ell+1]\sigma_{\ell+1}^A} \Lambda_R^{[\ell]}$, where $A$ is derived from $U$ and $\Lambda$ from $DV^\dagger$ in the way described before ($\Lambda_a B^\sigma_{ab} = M_{a\sigma,b} = \sum_k U_{a\sigma,k} D_k (V^\dagger)_{k,b} = A^\sigma_{ak} \Lambda_{kb}$). Similarly, we do a canonization from the right on  $A^{[\ell]\sigma_{\ell+1}^B} \Lambda^{[\ell]}$ to obtain $\Lambda_L^{[\ell]} B^{[\ell+1]\sigma_{\ell+1}^B}$, where $B$ is from $V^\dagger$. Then we have an ansatz
\begin{equation}
\ket{\psi_{\ell+1}^{{\rm guess}}} = \sum_{\fat{\sigma}} A^{[1]\sigma_1} \ldots  A^{[\ell+1]\sigma_{\ell+1}} \Lambda^{[\ell+1]}_{{\rm guess}} B^{[\ell+1]\sigma_{\ell+1}}  \ldots B^{[1]\sigma_1} \ket{\fat{\sigma}} ,
\end{equation}
where 
\begin{equation}
\Lambda^{[\ell+1]}_{{\rm guess}} =  \Lambda_R^{[\ell]} [\Lambda^{[\ell-1]}]^{-1} \Lambda^{[\ell]}_L .
\end{equation}
From this ansatz, we can then iteratively find the $\Lambda^{[\ell+1]}$ that minimizes the energy, slightly modifying the minimization part of variational MPS for a single site. In general, the result will not have the diagonal form resulting from an SVD, because $\Lambda_R$ and $\Lambda_L$ are not diagonal to begin with. But an SVD on it yields two unitaries that can be absorbed into the neighbouring $A$ and $B$ without affecting their normalization properties, such that the final $\Lambda^{[\ell+1]}$ is diagonal. In this form, the algorithm can be represented as in Fig.~\ref{fig:infinitebuildup}.

\begin{figure}
\centering\includegraphics[scale=0.7]{PyMPS_InfiniteSystemDMRG.eps}
\caption{Summarizing representation of the infinite system algorithm including prediction. At each step, two more sites are introduced, with an ansatz (white) calculated from the last iteration, leading to a new singular value decomposition (black).}
\label{fig:infinitebuildup}
\end{figure}


As shown by McCulloch\cite{McCulloch08}, this prediction leads to a dramatic speedup of infinite-system DMRG which complements nicely prediction algorithms of finite-system DMRG: the overlap between the predicted and calculated state often approaches unity up to $10^{-10}$ or so!

\subsection{iDMRG: variational ground state search in the thermodynamic limit}
Using the ideas of the preceding sections, it is very simple now to turn infinite-system DMRG into a performing and precise algorithm, called {\em iDMRG}, referring to the thermodynamic limit version of DMRG:
\begin{itemize}
\item Set up an infinite-system DMRG algorithm.
\item Add the prediction procedure of the last section to the minimization algorithm.
\item Run the modified algorithm until convergence is achieved. Convergence of the wave function
to the thermodynamic limit can be judged by considering the relationship Eq.~(\ref{eq:rhoarecursion})
\begin{equation}
\sum_{\sigma_\ell} A^{[\ell]\sigma_\ell} \rho_A^{[\ell]} A^{[\ell]\sigma_\ell\dagger} = \rho_A^{[\ell-1]} ,
\end{equation}
where $\rho_A^{[\ell]} = \Lambda^{[\ell]}\Lambda^{[\ell]\dagger}$ and $\rho_A^{[\ell-1]} = \Lambda^{[\ell-1]}\Lambda^{[\ell-1]\dagger}$ are the reduced density operators of the left half of the system. Note that while this relationship holds always between reduced density operators in the same finite system, here they originate from systems of two different lengths $2(\ell-1)$ and $2\ell$, such that this relationship is expected to hold only as a fixed point relationship for $\ell\rightarrow\infty$. Following the same argument as for generating the ansatz for the larger system, we may transform $A^{[\ell]\sigma_\ell}\Lambda^{[\ell]}$ to $\Lambda^{[\ell]}_L B^{[\ell+1]\sigma_\ell}$. Then the left-hand side of the fixed point relationship simplifies, using right normalization, to $\Lambda^{[\ell]}_L\Lambda^{[\ell]\dagger}_L \equiv \hat{\rho}^{[\ell]}_L$, and it becomes $\hat{\rho}^{[\ell]}_L = \rho_A^{[\ell-1]}$. If this relationship holds to high accuracy, the thermodynamic fixed point has been reached.  
One way of measuring the closeness of the two density operators is given by the fidelity \cite{McCulloch08}
\begin{equation}
F(\hat{\rho}^{[\ell]}_L, \hat{\rho}_A^{[\ell-1]}) = \tr \sqrt{\sqrt{\hat{\rho}^{[\ell]}_L}\hat{\rho}_A^{[\ell-1]} \sqrt{\hat{\rho}^{[\ell]}_L}} .
\end{equation}
Inserting the definitions and using cyclicity properties of the trace, one can show that $F(\hat{\rho}^{[\ell]}_L, \hat{\rho}^{[\ell-1]}_A) = \sum_i s_i$, where $s_i$ are the singular values of  $\Lambda^{[\ell]}_L \Lambda^{[\ell-1]\dagger}$. 
\end{itemize}
Of course the algorithm can be stopped at any time, but then we have a finite system result which definitely can be improved by using finite-system DMRG. The convergence criterion given really gives us access to the thermodynamic limit state, which we might write down formally as
\begin{equation}
\ket{\psi} = \sum_{\fat{\sigma}} \ldots A^{[\ell]\sigma_i}\Lambda^{[\ell]} B^{[\ell]\sigma_{i+1}} [\Lambda^{[\ell-1]}]^{-1}A^{[\ell]\sigma_{i+2}}\Lambda^{[\ell]} B^{[\ell]\sigma_{i+3}} [\Lambda^{[\ell-1]}]^{-1}A^{[\ell]\sigma_{i+4}}\Lambda^{[\ell]} B^{[\ell]\sigma_{i+5}} [\Lambda^{[\ell-1]}]^{-1} \ldots 
 \ket{\fat{\sigma}} ,
\end{equation}
where we take $\ell$ to be the iteration step when the convergence criterion is met. The question is now how to evaluate expectation values. Obviously, we cannot write down a finite network contraction as before; it will be of infinite size and therefore cannot be contracted naively. A contraction can only work if we can reduce the number to a finite number of contractions. For finite-system networks we saw that left- and right-orthonormality allow to eliminate most contractions: For observables on sites $i$ and $j$, one only had to consider the contractions on and between these two sites. There is however no reason why the thermodynamic limit state should meet normalization criteria; in fact, usually it does not. We therefore need an orthonormalization procedure. After that, expectation values can be evaluated as for a finite lattice with left- and right-orthonormalization. Because this procedure is also important for the next algorithm and conceptually a bit more advanced, I postpone it to an extra section. 

